% Options for packages loaded elsewhere
\PassOptionsToPackage{unicode}{hyperref}
\PassOptionsToPackage{hyphens}{url}
%
\documentclass[
]{article}
\usepackage{amsmath,amssymb}
\usepackage{iftex}
\ifPDFTeX
  \usepackage[T1]{fontenc}
  \usepackage[utf8]{inputenc}
  \usepackage{textcomp} % provide euro and other symbols
\else % if luatex or xetex
  \usepackage{unicode-math} % this also loads fontspec
  \defaultfontfeatures{Scale=MatchLowercase}
  \defaultfontfeatures[\rmfamily]{Ligatures=TeX,Scale=1}
\fi
\usepackage{lmodern}
\ifPDFTeX\else
  % xetex/luatex font selection
\fi
% Use upquote if available, for straight quotes in verbatim environments
\IfFileExists{upquote.sty}{\usepackage{upquote}}{}
\IfFileExists{microtype.sty}{% use microtype if available
  \usepackage[]{microtype}
  \UseMicrotypeSet[protrusion]{basicmath} % disable protrusion for tt fonts
}{}
\makeatletter
\@ifundefined{KOMAClassName}{% if non-KOMA class
  \IfFileExists{parskip.sty}{%
    \usepackage{parskip}
  }{% else
    \setlength{\parindent}{0pt}
    \setlength{\parskip}{6pt plus 2pt minus 1pt}}
}{% if KOMA class
  \KOMAoptions{parskip=half}}
\makeatother
\usepackage{xcolor}
\usepackage[margin=1in]{geometry}
\usepackage{color}
\usepackage{fancyvrb}
\newcommand{\VerbBar}{|}
\newcommand{\VERB}{\Verb[commandchars=\\\{\}]}
\DefineVerbatimEnvironment{Highlighting}{Verbatim}{commandchars=\\\{\}}
% Add ',fontsize=\small' for more characters per line
\usepackage{framed}
\definecolor{shadecolor}{RGB}{248,248,248}
\newenvironment{Shaded}{\begin{snugshade}}{\end{snugshade}}
\newcommand{\AlertTok}[1]{\textcolor[rgb]{0.94,0.16,0.16}{#1}}
\newcommand{\AnnotationTok}[1]{\textcolor[rgb]{0.56,0.35,0.01}{\textbf{\textit{#1}}}}
\newcommand{\AttributeTok}[1]{\textcolor[rgb]{0.13,0.29,0.53}{#1}}
\newcommand{\BaseNTok}[1]{\textcolor[rgb]{0.00,0.00,0.81}{#1}}
\newcommand{\BuiltInTok}[1]{#1}
\newcommand{\CharTok}[1]{\textcolor[rgb]{0.31,0.60,0.02}{#1}}
\newcommand{\CommentTok}[1]{\textcolor[rgb]{0.56,0.35,0.01}{\textit{#1}}}
\newcommand{\CommentVarTok}[1]{\textcolor[rgb]{0.56,0.35,0.01}{\textbf{\textit{#1}}}}
\newcommand{\ConstantTok}[1]{\textcolor[rgb]{0.56,0.35,0.01}{#1}}
\newcommand{\ControlFlowTok}[1]{\textcolor[rgb]{0.13,0.29,0.53}{\textbf{#1}}}
\newcommand{\DataTypeTok}[1]{\textcolor[rgb]{0.13,0.29,0.53}{#1}}
\newcommand{\DecValTok}[1]{\textcolor[rgb]{0.00,0.00,0.81}{#1}}
\newcommand{\DocumentationTok}[1]{\textcolor[rgb]{0.56,0.35,0.01}{\textbf{\textit{#1}}}}
\newcommand{\ErrorTok}[1]{\textcolor[rgb]{0.64,0.00,0.00}{\textbf{#1}}}
\newcommand{\ExtensionTok}[1]{#1}
\newcommand{\FloatTok}[1]{\textcolor[rgb]{0.00,0.00,0.81}{#1}}
\newcommand{\FunctionTok}[1]{\textcolor[rgb]{0.13,0.29,0.53}{\textbf{#1}}}
\newcommand{\ImportTok}[1]{#1}
\newcommand{\InformationTok}[1]{\textcolor[rgb]{0.56,0.35,0.01}{\textbf{\textit{#1}}}}
\newcommand{\KeywordTok}[1]{\textcolor[rgb]{0.13,0.29,0.53}{\textbf{#1}}}
\newcommand{\NormalTok}[1]{#1}
\newcommand{\OperatorTok}[1]{\textcolor[rgb]{0.81,0.36,0.00}{\textbf{#1}}}
\newcommand{\OtherTok}[1]{\textcolor[rgb]{0.56,0.35,0.01}{#1}}
\newcommand{\PreprocessorTok}[1]{\textcolor[rgb]{0.56,0.35,0.01}{\textit{#1}}}
\newcommand{\RegionMarkerTok}[1]{#1}
\newcommand{\SpecialCharTok}[1]{\textcolor[rgb]{0.81,0.36,0.00}{\textbf{#1}}}
\newcommand{\SpecialStringTok}[1]{\textcolor[rgb]{0.31,0.60,0.02}{#1}}
\newcommand{\StringTok}[1]{\textcolor[rgb]{0.31,0.60,0.02}{#1}}
\newcommand{\VariableTok}[1]{\textcolor[rgb]{0.00,0.00,0.00}{#1}}
\newcommand{\VerbatimStringTok}[1]{\textcolor[rgb]{0.31,0.60,0.02}{#1}}
\newcommand{\WarningTok}[1]{\textcolor[rgb]{0.56,0.35,0.01}{\textbf{\textit{#1}}}}
\usepackage{graphicx}
\makeatletter
\newsavebox\pandoc@box
\newcommand*\pandocbounded[1]{% scales image to fit in text height/width
  \sbox\pandoc@box{#1}%
  \Gscale@div\@tempa{\textheight}{\dimexpr\ht\pandoc@box+\dp\pandoc@box\relax}%
  \Gscale@div\@tempb{\linewidth}{\wd\pandoc@box}%
  \ifdim\@tempb\p@<\@tempa\p@\let\@tempa\@tempb\fi% select the smaller of both
  \ifdim\@tempa\p@<\p@\scalebox{\@tempa}{\usebox\pandoc@box}%
  \else\usebox{\pandoc@box}%
  \fi%
}
% Set default figure placement to htbp
\def\fps@figure{htbp}
\makeatother
\setlength{\emergencystretch}{3em} % prevent overfull lines
\providecommand{\tightlist}{%
  \setlength{\itemsep}{0pt}\setlength{\parskip}{0pt}}
\setcounter{secnumdepth}{-\maxdimen} % remove section numbering
\usepackage{bookmark}
\IfFileExists{xurl.sty}{\usepackage{xurl}}{} % add URL line breaks if available
\urlstyle{same}
\hypersetup{
  pdftitle={STAT 325 - Design of Experiment Homework 3},
  pdfauthor={Ducheng Tan},
  hidelinks,
  pdfcreator={LaTeX via pandoc}}

\title{STAT 325 - Design of Experiment Homework 3}
\author{Ducheng Tan}
\date{2026-02-1}

\begin{document}
\maketitle

\textbf{Question 1}

P-value calculation

\begin{Shaded}
\begin{Highlighting}[]
\NormalTok{p\_value }\OtherTok{\textless{}{-}} \DecValTok{1} \SpecialCharTok{{-}} \FunctionTok{pf}\NormalTok{(}\FloatTok{4.3063}\NormalTok{, }\AttributeTok{df1 =} \DecValTok{4}\NormalTok{, }\AttributeTok{df2 =} \DecValTok{60}\NormalTok{)}
\NormalTok{p\_value}
\end{Highlighting}
\end{Shaded}

\begin{verbatim}
## [1] 0.003960773
\end{verbatim}

\textbf{Question 3}

\begin{Shaded}
\begin{Highlighting}[]
\NormalTok{data }\OtherTok{\textless{}{-}} \FunctionTok{data.frame}\NormalTok{(}\AttributeTok{diet =} \FunctionTok{c}\NormalTok{(}\DecValTok{1}\NormalTok{, }\DecValTok{1}\NormalTok{, }\DecValTok{1}\NormalTok{,}
           \DecValTok{2}\NormalTok{, }\DecValTok{2}\NormalTok{, }\DecValTok{2}\NormalTok{, }\DecValTok{2}\NormalTok{, }\DecValTok{2}\NormalTok{,}
           \DecValTok{3}\NormalTok{, }\DecValTok{3}\NormalTok{, }\DecValTok{3}\NormalTok{, }\DecValTok{3}\NormalTok{), }
           
\AttributeTok{WeightGain =} \FunctionTok{c}\NormalTok{(}\DecValTok{8}\NormalTok{, }\DecValTok{16}\NormalTok{, }\DecValTok{9}\NormalTok{,}
                 \DecValTok{9}\NormalTok{, }\DecValTok{16}\NormalTok{, }\DecValTok{21}\NormalTok{, }\DecValTok{11}\NormalTok{, }\DecValTok{18}\NormalTok{,}
                 \DecValTok{15}\NormalTok{, }\DecValTok{10}\NormalTok{, }\DecValTok{17}\NormalTok{, }\DecValTok{6}\NormalTok{)}
\NormalTok{)}


\NormalTok{data\_2}\OtherTok{\textless{}{-}} \FunctionTok{data.frame}\NormalTok{(}
  \AttributeTok{diet =} \FunctionTok{factor}\NormalTok{(}\FunctionTok{c}\NormalTok{(}
    \FunctionTok{rep}\NormalTok{(}\StringTok{"1"}\NormalTok{, }\DecValTok{3}\NormalTok{),}
    \FunctionTok{rep}\NormalTok{(}\StringTok{"2"}\NormalTok{, }\DecValTok{5}\NormalTok{),}
    \FunctionTok{rep}\NormalTok{(}\StringTok{"3"}\NormalTok{, }\DecValTok{4}\NormalTok{)}
\NormalTok{  )),}
  \AttributeTok{WeightGain =} \FunctionTok{c}\NormalTok{(}
    \DecValTok{8}\NormalTok{, }\DecValTok{16}\NormalTok{, }\DecValTok{9}\NormalTok{,}
    \DecValTok{9}\NormalTok{, }\DecValTok{16}\NormalTok{, }\DecValTok{21}\NormalTok{, }\DecValTok{11}\NormalTok{, }\DecValTok{18}\NormalTok{,}
    \DecValTok{15}\NormalTok{, }\DecValTok{10}\NormalTok{, }\DecValTok{17}\NormalTok{, }\DecValTok{6}
\NormalTok{  )}
\NormalTok{)}
\FunctionTok{head}\NormalTok{(data\_2)}
\end{Highlighting}
\end{Shaded}

\begin{verbatim}
##   diet WeightGain
## 1    1          8
## 2    1         16
## 3    1          9
## 4    2          9
## 5    2         16
## 6    2         21
\end{verbatim}

\begin{Shaded}
\begin{Highlighting}[]
\FunctionTok{str}\NormalTok{(data)}
\end{Highlighting}
\end{Shaded}

\begin{verbatim}
## 'data.frame':    12 obs. of  2 variables:
##  $ diet      : num  1 1 1 2 2 2 2 2 3 3 ...
##  $ WeightGain: num  8 16 9 9 16 21 11 18 15 10 ...
\end{verbatim}

\begin{Shaded}
\begin{Highlighting}[]
\FunctionTok{head}\NormalTok{(data)}
\end{Highlighting}
\end{Shaded}

\begin{verbatim}
##   diet WeightGain
## 1    1          8
## 2    1         16
## 3    1          9
## 4    2          9
## 5    2         16
## 6    2         21
\end{verbatim}

\begin{Shaded}
\begin{Highlighting}[]
\CommentTok{\#View(data)}
\end{Highlighting}
\end{Shaded}

\begin{Shaded}
\begin{Highlighting}[]
\FunctionTok{table}\NormalTok{(data}\SpecialCharTok{$}\NormalTok{diet)}
\end{Highlighting}
\end{Shaded}

\begin{verbatim}
## 
## 1 2 3 
## 3 5 4
\end{verbatim}

\begin{Shaded}
\begin{Highlighting}[]
\FunctionTok{summary}\NormalTok{(data)}
\end{Highlighting}
\end{Shaded}

\begin{verbatim}
##       diet         WeightGain   
##  Min.   :1.000   Min.   : 6.00  
##  1st Qu.:1.750   1st Qu.: 9.00  
##  Median :2.000   Median :13.00  
##  Mean   :2.083   Mean   :13.00  
##  3rd Qu.:3.000   3rd Qu.:16.25  
##  Max.   :3.000   Max.   :21.00
\end{verbatim}

\begin{Shaded}
\begin{Highlighting}[]
\FunctionTok{tapply}\NormalTok{(data}\SpecialCharTok{$}\NormalTok{WeightGain, data}\SpecialCharTok{$}\NormalTok{diet, summary)}
\end{Highlighting}
\end{Shaded}

\begin{verbatim}
## $`1`
##    Min. 1st Qu.  Median    Mean 3rd Qu.    Max. 
##     8.0     8.5     9.0    11.0    12.5    16.0 
## 
## $`2`
##    Min. 1st Qu.  Median    Mean 3rd Qu.    Max. 
##       9      11      16      15      18      21 
## 
## $`3`
##    Min. 1st Qu.  Median    Mean 3rd Qu.    Max. 
##     6.0     9.0    12.5    12.0    15.5    17.0
\end{verbatim}

\begin{Shaded}
\begin{Highlighting}[]
\FunctionTok{boxplot}\NormalTok{(WeightGain }\SpecialCharTok{\textasciitilde{}}\NormalTok{ diet, }\AttributeTok{data =}\NormalTok{ data,}
        \AttributeTok{xlab =} \StringTok{"Diet"}\NormalTok{,}
        \AttributeTok{ylab =} \StringTok{"Weight Gain"}\NormalTok{,}
        \AttributeTok{main =} \StringTok{"Lamb Weight Gain by Diet"}\NormalTok{)}
\end{Highlighting}
\end{Shaded}

\pandocbounded{\includegraphics[keepaspectratio]{HW3_files/figure-latex/unnamed-chunk-4-1.pdf}}
\textbf{How do the samples from the three diets compare?}

\textit{From basic descriptive statistics and the boxplot we can see that diet 2 corresponds with the most weight gain the this data by distribution and sample mean. Also Diet 3 is almost symmetric with a mean of 12. Lastly we can notice that diet 1 corresponds to the least weight gain with lower variation than the other two groups in this data.}

\textbf{Are the mean weight gains of the three diets the same?}

\textit{Although from basic descriptive summary statistics we can see that the mean weight gains of the three diets are not the same but that can be due to small sampling variation, we can also perform One-Way ANOVA to analyze the mean.}

\$\$ H\_0: \mu\_1 = \mu\_2= \mu\_3 \textbackslash{} H\_a: H\_0
\text{ is false }

\$\$
\textit{From the ANOVA table we can see that the P value $Pr(>F) = 0.491 >> \alpha = 0.05$ , this $P$ value is greater than any statistically significant level, thus we can say that there is no statistical evidence to reject $H_0$. In this context, there is insufficient statistical evidence to conclude that the mean weights gains in goats differ among the three diets.}

\begin{Shaded}
\begin{Highlighting}[]
\NormalTok{anova\_model }\OtherTok{\textless{}{-}} \FunctionTok{aov}\NormalTok{(WeightGain }\SpecialCharTok{\textasciitilde{}}\NormalTok{ diet, }\AttributeTok{data =}\NormalTok{ data\_2)}
\FunctionTok{summary}\NormalTok{(anova\_model)}
\end{Highlighting}
\end{Shaded}

\begin{verbatim}
##             Df Sum Sq Mean Sq F value Pr(>F)
## diet         2     36   18.00   0.771  0.491
## Residuals    9    210   23.33
\end{verbatim}

\textbf{If the means are not the same, which ones are different?} Since
the one-way ANOVA failed to reject the null hypothesis, there is no
statistical evidence that the mean weight gains differ among the three
diets

\textbf{Have the assumptions of the one-way ANOVA been met?}

In our model assumptions for using ANOVA we have

\begin{itemize}
\item - Independent Observations 
\item Errors are normally distributed 
\item Constant Variance 
\end{itemize}

For independent observations we can assume the errors are independent
across the lambs due to the design so
\[ \varepsilon_{ij} \neq \varepsilon_{i'j'}, \quad (i , j) \neq (i' ,j') \]
We check normality with histogram of residuals, Shapiro-Wilk Test, and
Normal probability plot/ QQ plot. With Shapiro-Wilk test we have the
hypothesis as

\[
 H_0: \text{Residuals are normally distributed} \\
 H_a : \text{Residuals are not normally distributed }
\]

\begin{Shaded}
\begin{Highlighting}[]
\FunctionTok{par}\NormalTok{(}\AttributeTok{mfrow =} \FunctionTok{c}\NormalTok{(}\DecValTok{2}\NormalTok{,}\DecValTok{2}\NormalTok{))}
\FunctionTok{plot}\NormalTok{(anova\_model)}
\end{Highlighting}
\end{Shaded}

\pandocbounded{\includegraphics[keepaspectratio]{HW3_files/figure-latex/unnamed-chunk-6-1.pdf}}

\begin{Shaded}
\begin{Highlighting}[]
\NormalTok{resi }\OtherTok{\textless{}{-}} \FunctionTok{residuals}\NormalTok{(anova\_model)}
\FunctionTok{shapiro.test}\NormalTok{(resi)}
\end{Highlighting}
\end{Shaded}

\begin{verbatim}
## 
##  Shapiro-Wilk normality test
## 
## data:  resi
## W = 0.91275, p-value = 0.2313
\end{verbatim}

\begin{Shaded}
\begin{Highlighting}[]
\FunctionTok{hist}\NormalTok{(resi)}
\end{Highlighting}
\end{Shaded}

\pandocbounded{\includegraphics[keepaspectratio]{HW3_files/figure-latex/unnamed-chunk-6-2.pdf}}
We comment that the residual histogram seems very normal, and the
QQ-plot of the residuals aligns with the Gaussian distribution. Then
lastly, we can see that a numeric test using Shapiro Wilk conveys that
we don't have significant evidence to reject \(H_0\) (The residuals are
normally distributed). \$P=0.5394 \textgreater\textgreater{} \alpha =
0.05 \$.

Next we test constant variance assumption with looking at the residual
vs fitted plot and using Bartlett's test and Levene' test.

\begin{Shaded}
\begin{Highlighting}[]
\FunctionTok{plot}\NormalTok{(}\FunctionTok{fitted}\NormalTok{(anova\_model), }\FunctionTok{residuals}\NormalTok{(anova\_model),}
     \AttributeTok{xlab =} \StringTok{"Fitted values"}\NormalTok{,}
     \AttributeTok{ylab =} \StringTok{"Residuals"}\NormalTok{,}
     \AttributeTok{main =} \StringTok{"Residuals vs Fitted Values"}\NormalTok{)}
\FunctionTok{abline}\NormalTok{(}\AttributeTok{h =} \DecValTok{0}\NormalTok{, }\AttributeTok{lty =} \DecValTok{2}\NormalTok{, }\AttributeTok{col =} \StringTok{"orange"}\NormalTok{)}
\end{Highlighting}
\end{Shaded}

\pandocbounded{\includegraphics[keepaspectratio]{HW3_files/figure-latex/unnamed-chunk-7-1.pdf}}
The plot of the residuals seems scattered in the same fashion thus we
can assume constant variance.

\begin{Shaded}
\begin{Highlighting}[]
\FunctionTok{bartlett.test}\NormalTok{(WeightGain }\SpecialCharTok{\textasciitilde{}}\NormalTok{ diet, }\AttributeTok{data =}\NormalTok{ data\_2)}
\end{Highlighting}
\end{Shaded}

\begin{verbatim}
## 
##  Bartlett test of homogeneity of variances
## 
## data:  WeightGain by diet
## Bartlett's K-squared = 0.042183, df = 2, p-value = 0.9791
\end{verbatim}

\begin{Shaded}
\begin{Highlighting}[]
\FunctionTok{library}\NormalTok{(car)}
\end{Highlighting}
\end{Shaded}

\begin{verbatim}
## Loading required package: carData
\end{verbatim}

\begin{Shaded}
\begin{Highlighting}[]
\FunctionTok{leveneTest}\NormalTok{(WeightGain }\SpecialCharTok{\textasciitilde{}}\NormalTok{ diet, }\AttributeTok{data =}\NormalTok{ data\_2)}
\end{Highlighting}
\end{Shaded}

\begin{verbatim}
## Levene's Test for Homogeneity of Variance (center = median)
##       Df F value Pr(>F)
## group  2  0.2203 0.8065
##        9
\end{verbatim}

From both test we see that we have a large p-value which means that we
fail to reject that the variances are equal. Lastly, potential outliers
were examined using standardized residuals and Cook's distance, and no
influential observations were identified.

\begin{Shaded}
\begin{Highlighting}[]
\FunctionTok{plot}\NormalTok{(}\FunctionTok{cooks.distance}\NormalTok{(anova\_model),}
     \AttributeTok{type =} \StringTok{"h"}\NormalTok{,}
     \AttributeTok{ylab =} \StringTok{"Cook\textquotesingle{}s Distance"}\NormalTok{,}
     \AttributeTok{main =} \StringTok{"Cook\textquotesingle{}s Distance"}\NormalTok{)}
\FunctionTok{abline}\NormalTok{(}\AttributeTok{h =} \DecValTok{4}\SpecialCharTok{/}\FunctionTok{length}\NormalTok{(}\FunctionTok{residuals}\NormalTok{(anova\_model)), }\AttributeTok{lty =} \DecValTok{2}\NormalTok{)}
\end{Highlighting}
\end{Shaded}

\pandocbounded{\includegraphics[keepaspectratio]{HW3_files/figure-latex/unnamed-chunk-9-1.pdf}}

\end{document}
